\section{Conclusion}
\label{conclusion}

\textbf{Answer.} On EMSCAD, under one protocol with grouped splits and five seeds, a logistic
regression on BM25-weighted words is the best single model (test PR-AUC 0.847), and fine-tuned
DistilBERT does not beat it (0.816). Stacking the two gives the best result, 0.878, and the gain over
BM25 LR is positive on all five seeds. Detectors that never see a fraud label (one-class SVM,
zero-shot Qwen2.5-1.5B) stay near the 0.049 floor.

\textbf{Lessons.}
\begin{enumerate}[nosep]
  \item Split by description group, not by row: EMSCAD repeats descriptions, and a row split would
        reward memorising them.
  \item Report PR-AUC and the spread over several splits; compare models paired by split. The
        unpaired error bars hide consistent gains such as BM25 over TF-IDF (5 of 5 seeds).
  \item Start from a tuned linear baseline. It was the strongest single model here at a fraction of
        the cost of the alternatives.
  \item Two models with similar scores but different mistakes are worth combining; a two-input
        stacker on validation scores is cheap and gave the best result of the project.
  \item Read the explanations for artefacts: the strongest features include two countries, one
        employer name and one URL hash, which say more about the dataset than about fraud.
\end{enumerate}

\textbf{Future work.} Check the country and template artefacts by training without the country
feature and with splits grouped by poster rather than by description; try the stack with the
MLP as a third input; and test a model trained on EMSCAD on job ads from a different source.
